\chapter*{\centering ABSTRAK}
\addcontentsline{toc}{chapter}{ABSTRAK}

\begin{spacing}{1.25}
\noindent Abstrak memberikan gambaran singkat terhadap usulan penelitian dan memuat secara ringkas: (1) latar belakang permasalahan, (2) tujuan penelitian, (3) urgensi usulan penelitian, yaitu alasan mengapa permasalahan tersebut penting dan mendesak untuk diselesaikan serta \textit{research gap} yang hendak dijawab, dan (4) metode penelitian yang diajukan. Karena merupakan usulan, abstrak proposal tidak memuat hasil penelitian. Panjang abstrak antara 150--250 kata, ditulis dalam \textit{satu paragraf}, dan tidak mengandung rujukan. Hindari penggunaan singkatan yang tidak didefinisikan, persamaan matematika, atau tabel.\par\medskip

\noindent\textbf{Kata kunci:} [Kata kunci 1], [Kata kunci 2], [Kata kunci 3], [Kata kunci 4], [Kata kunci 5] \textit{(pilih 3--5 kata/frasa yang paling mencerminkan topik penelitian, dipisahkan oleh koma)}
\end{spacing}
